\scanpage{183}
\begin{center}- 31 -\end{center}\hfill IX

3.3.2. On voit donc que les propriétés équivalentes de 3.2 impliquent
une propriété de ``stabilité'' partielle pour les modèles de Néron ; on
fera attention que lorsque $A_K$ n'a pas bonne réduction i.e. que $A$ n'est
pas un schéma abélien, il n'est \uline{pas} vrai que pour tout $S'$ comme dans
3.3 (et même pour $S'$ fini sur $S$, et $K'/K$ séparable $\cdots$), le morphisme
$A\times_S S'\longrightarrow A'$ lui-même soit un isomorphisme, i.e. induise des isomorphismes
sur les groupes de composantes connexes des fibres (cf. 11.10 plus bas
pour le comportement de ces groupes de composantes connexes par extension
de la base). Ces observations, renforcées par 3.9 plus bas, justifient
la terminologie suivante, introduite par MUMFORD (et qui s'harmonise
aussi, paraît-il, avec les notions de même nom étudiées dans son livre [17]):

\uline{Définition} 3.4. \uline{Les notations étant celles de} 3.2, \uline{on dit que le schéma
abélien} $A_K$ \uline{sur} $K$ \uline{a} réduction semi-stable sur $S$ \uline{si les conditions équivalentes} (i) \uline{et} (ii) \uline{de} 3.2 \uline{sont vérifiées}.

3.4.0. On notera que sous la forme (ii), cette condition ne fait pas
appel à la notion de modèle de Néron. D'autre part, elle se ramène
immédiatement par localisation au cas où $S$ est un trait, comme il
résulte par exemple de 3.2.1. D'ailleurs, comme la formation du modèle
de Néron commute au changement de base de $S$ vers son hensélisé, ou
son hensélisé strict, la notion se ramène même au cas $S$ où $S$ est
un trait strictement local.
\margin{\unsure{345|349}}
\margin{349}
\margin{182}
